\section{The definition of a triangulated category}
\label{section-triangulated-definitions}

\noindent
In this section we collect most of the definitions concerning triangulated
and pre-triangulated categories.

\begin{definition}
\label{definition-triangle}
Let $\mathcal{D}$ be an additive category. Let
$[1] : \mathcal{D} \to \mathcal{D}$, $E \mapsto E[1]$
be an additive functor which is an auto-equivalence of $\mathcal{D}$.
\begin{enumerate}
\item A {\it triangle} is a sextuple
$(X, Y, Z, f, g, h)$ where $X, Y, Z \in \Ob(\mathcal{D})$ and
$f : X \to Y$, $g : Y \to Z$ and $h : Z \to X[1]$ are morphisms
of $\mathcal{D}$.
\item A {\it morphism of triangles}
$(X, Y, Z, f, g, h) \to (X', Y', Z', f', g', h')$
is given by morphisms $a : X \to X'$, $b : Y \to Y'$ and $c : Z \to Z'$
of $\mathcal{D}$ such that
$b \circ f = f' \circ a$, $c  \circ g = g' \circ b$ and
$a[1] \circ h = h' \circ c$.
\end{enumerate}
\end{definition}

\noindent
A morphism of triangles is visualized by the following
commutative diagram
$$
\xymatrix{
X \ar[r] \ar[d]^a &
Y \ar[r] \ar[d]^b &
Z \ar[r] \ar[d]^c &
X[1] \ar[d]^{a[1]} \\
X' \ar[r] &
Y' \ar[r] &
Z' \ar[r] &
X'[1]
}
$$
In the setting of Definition \ref{definition-triangle}, we
write $[0] = \text{id}$, for $n > 0$ we denote $[n]$ the $n$-fold
composition of $[1]$, we choose a quasi-inverse $[-1]$ of $[1]$, and
we set $[-n]$ equal to the $n$-fold composition of $[-1]$. Then
$\{[n]\}_{n \in \mathbf{Z}}$ is a collection of additive auto-equivalences
of $\mathcal{D}$ indexed by $n \in \mathbf{Z}$ such that we are given
isomorphisms of functors $[n] \circ [m] \cong [n + m]$.

\medskip\noindent
Here is the definition of a triangulated category as given in
Verdier's thesis.

\begin{definition}
\label{definition-triangulated-category}
A {\it triangulated category} consists of a triple
$(\mathcal{D}, \{[n]\}_{n\in \mathbf{Z}}, \mathcal{T})$
where
\begin{enumerate}
\item $\mathcal{D}$ is an additive category,
\item $[1] : \mathcal{D} \to \mathcal{D}$, $E \mapsto E[1]$
is an additive auto-equivalence and $[n]$ for $n \in \mathbf{Z}$ is
as discussed above, and
\item $\mathcal{T}$ is a set of triangles (Definition \ref{definition-triangle})
called the {\it distinguished triangles}
\end{enumerate}
subject to the following conditions
\begin{enumerate}
\item[TR1] Any triangle isomorphic to a distinguished triangle is
a distinguished triangle. Any triangle of the form
$(X, X, 0, \text{id}, 0, 0)$ is distinguished.
For any morphism $f : X \to Y$ of $\mathcal{D}$ there exists a
distinguished triangle of the form $(X, Y, Z, f, g, h)$.
\item[TR2] The triangle $(X, Y, Z, f, g, h)$ is distinguished
if and only if the triangle $(Y, Z, X[1], g, h, -f[1])$ is.
\item[TR3] Given a solid diagram
$$
\xymatrix{
X \ar[r]^f \ar[d]^a &
Y \ar[r]^g \ar[d]^b &
Z \ar[r]^h \ar@{-->}[d] &
X[1] \ar[d]^{a[1]} \\
X' \ar[r]^{f'} &
Y' \ar[r]^{g'} &
Z' \ar[r]^{h'} &
X'[1]
}
$$
whose rows are distinguished triangles and which satisfies
$b \circ f = f' \circ a$, there exists a morphism
$c : Z \to Z'$ such that $(a, b, c)$ is a morphism of triangles.
\item[TR4] Given objects $X$, $Y$, $Z$ of $\mathcal{D}$, and morphisms
$f : X \to Y$, $g : Y \to Z$, and distinguished triangles
$(X, Y, Q_1, f, p_1, d_1)$,
$(X, Z, Q_2, g \circ f, p_2, d_2)$,
and
$(Y, Z, Q_3, g, p_3, d_3)$,
there exist
morphisms $a : Q_1 \to Q_2$ and $b : Q_2 \to Q_3$ such
that
\begin{enumerate}
\item $(Q_1, Q_2, Q_3, a, b, p_1[1] \circ d_3)$ is a
distinguished triangle,
\item the triple $(\text{id}_X, g, a)$ is
a morphism of triangles
$(X, Y, Q_1, f, p_1, d_1) \to (X, Z, Q_2, g \circ f, p_2, d_2)$, and
\item the triple $(f, \text{id}_Z, b)$ is a morphism of triangles
$(X, Z, Q_2, g \circ f, p_2, d_2) \to (Y, Z, Q_3, g, p_3, d_3)$.
\end{enumerate}
\end{enumerate}
We will call $(\mathcal{D}, [\ ], \mathcal{T})$ a
{\it pre-triangulated category} if TR1, TR2 and TR3
hold.\footnote{We use $[\ ]$ as an abbreviation for the
family $\{[n]\}_{n\in \mathbf{Z}}$.}
\end{definition}

\noindent
The explanation of TR4 is that if you think of $Q_1$ as
$Y/X$, $Q_2$ as $Z/X$ and $Q_3$ as $Z/Y$, then TR4(a) expresses
the isomorphism $(Z/X)/(Y/X) \cong Z/Y$ and TR4(b) and TR4(c)
express that we can compare the triangles $X \to Y \to Q_1 \to X[1]$
etc with morphisms of triangles. For a more precise reformulation
of this idea see the proof of Lemma \ref{lemma-two-split-injections}.

\medskip\noindent
The sign in TR2 means that if $(X, Y, Z, f, g, h)$ is a distinguished triangle
then in the long sequence
\begin{equation}
\label{equation-rotate}
\ldots \to
Z[-1] \xrightarrow{-h[-1]}
X \xrightarrow{f}
Y \xrightarrow{g}
Z \xrightarrow{h}
X[1] \xrightarrow{-f[1]}
Y[1] \xrightarrow{-g[1]}
Z[1] \to \ldots
\end{equation}
each four term sequence gives a distinguished triangle.

\medskip\noindent
As usual we abuse notation and we simply speak of a (pre-)triangulated
category $\mathcal{D}$ without explicitly introducing notation for the
additional data. The notion of a pre-triangulated category is
useful in finding statements equivalent to TR4.

\medskip\noindent
We have the following definition of a triangulated functor.

\begin{definition}
\label{definition-exact-functor-triangulated-categories}
Let $\mathcal{D}$, $\mathcal{D}'$ be pre-triangulated
categories. An {\it exact functor}, or a {\it triangulated functor}
from $\mathcal{D}$ to $\mathcal{D}'$ is a functor
$F : \mathcal{D} \to \mathcal{D}'$ together
with given functorial isomorphisms $\xi_X : F(X[1]) \to F(X)[1]$
such that for every distinguished triangle
$(X, Y, Z, f, g, h)$ of $\mathcal{D}$ the triangle
$(F(X), F(Y), F(Z), F(f), F(g), \xi_X \circ F(h))$
is a distinguished triangle of $\mathcal{D}'$.
\end{definition}

\noindent
An exact functor is additive, see
Lemma \ref{lemma-exact-functor-additive}.
When we say two triangulated categories are equivalent we mean that
they are equivalent in the $2$-category of triangulated categories.
A $2$-morphism $a : (F, \xi) \to (F', \xi')$ in this $2$-category is
simply a transformation of functors $a : F \to F'$ which is compatible
with $\xi$ and $\xi'$, i.e.,
$$
\xymatrix{
F \circ [1] \ar[r]_\xi \ar[d]_{a \star 1} & [1] \circ F \ar[d]^{1 \star a} \\
F' \circ [1] \ar[r]^{\xi'} & [1] \circ F'
}
$$
commutes.

\begin{definition}
\label{definition-triangulated-subcategory}
Let $(\mathcal{D}, [\ ], \mathcal{T})$ be a pre-triangulated category.
A {\it pre-triangulated subcategory}\footnote{This definition may be
nonstandard. If $\mathcal{D}'$ is a full subcategory then $\mathcal{T}'$
is the intersection of the set of triangles in $\mathcal{D}'$ with
$\mathcal{T}$, see
Lemma \ref{lemma-triangulated-subcategory}.
In this case we drop $\mathcal{T}'$ from the notation.}
is a pair $(\mathcal{D}', \mathcal{T}')$ such that
\begin{enumerate}
\item $\mathcal{D}'$ is an additive subcategory of $\mathcal{D}$
which is preserved under $[1]$ and such that
$[1] : \mathcal{D}' \to \mathcal{D}'$ is an auto-equivalence,
\item $\mathcal{T}' \subset \mathcal{T}$ is a subset such that for every
$(X, Y, Z, f, g, h) \in \mathcal{T}'$ we have
$X, Y, Z \in \Ob(\mathcal{D}')$ and
$f, g, h \in \text{Arrows}(\mathcal{D}')$, and
\item $(\mathcal{D}', [\ ], \mathcal{T}')$ is a pre-triangulated
category.
\end{enumerate}
If $\mathcal{D}$ is a triangulated category, then we say
$(\mathcal{D}', \mathcal{T}')$ is a {\it triangulated subcategory} if
it is a pre-triangulated subcategory and
$(\mathcal{D}', [\ ], \mathcal{T}')$ is a triangulated category.
\end{definition}

\noindent
In this situation the inclusion functor
$\mathcal{D}' \to \mathcal{D}$ is an exact functor
with $\xi_X : X[1] \to X[1]$ given by the identity on $X[1]$.

\medskip\noindent
We will see in
Lemma \ref{lemma-composition-zero}
that for a distinguished triangle $(X, Y, Z, f, g, h)$
in a pre-triangulated category the composition $g \circ f : X \to Z$ is zero.
Thus the sequence (\ref{equation-rotate}) is a complex.
A homological functor is one that turns this complex into a long
exact sequence.

\begin{definition}
\label{definition-homological}
Let $\mathcal{D}$ be a pre-triangulated category.
Let $\mathcal{A}$ be an abelian category.
An additive functor $H : \mathcal{D} \to \mathcal{A}$ is called
{\it homological} if for every distinguished triangle
$(X, Y, Z, f, g, h)$ the sequence
$$
H(X) \to H(Y) \to H(Z)
$$
is exact in the abelian category $\mathcal{A}$. An additive functor
$H : \mathcal{D}^{opp} \to \mathcal{A}$ is called {\it cohomological}
if the corresponding functor $\mathcal{D} \to \mathcal{A}^{opp}$ is
homological.
\end{definition}

\noindent
If $H : \mathcal{D} \to \mathcal{A}$ is a homological functor
we often write $H^n(X) = H(X[n])$ so that $H(X) = H^0(X)$.
Our discussion of TR2 above implies that a distinguished triangle
$(X, Y, Z, f, g, h)$ determines a long exact sequence
\begin{equation}
\label{equation-long-exact-cohomology-sequence}
\xymatrix@C=3pc{
H^{-1}(Z) \ar[r]^{H(h[-1])} &
H^0(X) \ar[r]^{H(f)} &
H^0(Y) \ar[r]^{H(g)} &
H^0(Z) \ar[r]^{H(h)} &
H^1(X)
}
\end{equation}
This will be called the {\it long exact sequence} associated to the
distinguished triangle and the homological functor. As indicated
we will not use any signs for the morphisms in the long exact
sequence. This has the side effect that maps in the long exact sequence
associated to the rotation (TR2) of a distinguished triangle differ
from the maps in the sequence above by some signs.

\begin{definition}
\label{definition-delta-functor}
Let $\mathcal{A}$ be an abelian category.
Let $\mathcal{D}$ be a triangulated category.
A {\it $\delta$-functor from $\mathcal{A}$ to $\mathcal{D}$} is
given by a functor $G : \mathcal{A} \to \mathcal{D}$ and
a rule which assigns to every short exact sequence
$$
0 \to A \xrightarrow{a} B \xrightarrow{b} C \to 0
$$
a morphism $\delta = \delta_{A \to B \to C} : G(C) \to G(A)[1]$
such that
\begin{enumerate}
\item the triangle
$(G(A), G(B), G(C), G(a), G(b), \delta_{A \to B \to C})$
is a distinguished triangle of $\mathcal{D}$
for any short exact sequence as above, and
\item for every morphism $(A \to B \to C) \to (A' \to B' \to C')$
of short exact sequences the diagram
$$
\xymatrix{
G(C) \ar[d] \ar[rr]_{\delta_{A \to B \to C}} & &
G(A)[1] \ar[d] \\
G(C') \ar[rr]^{\delta_{A' \to B' \to C'}} & &
G(A')[1]
}
$$
is commutative.
\end{enumerate}
In this situation we call
$(G(A), G(B), G(C), G(a), G(b), \delta_{A \to B \to C})$
the {\it image of the short exact sequence under the
given $\delta$-functor}.
\end{definition}

\noindent
Note how a $\delta$-functor comes equipped with additional structure.
Strictly speaking it does not make sense to say that a given
functor $\mathcal{A} \to \mathcal{D}$ is a $\delta$-functor, but we
will often do so anyway.












